\section*{Covariances and conventions: a reading guide}\addcontentsline{toc}{section}{Covariances
and conventions: a reading guide}

\noindent\revs{Texts on fermionic Gaussian states, for example~\cite{bravyi,ksvc,surace}, describe a
state of $n$ modes by a matrix of second moments of size $2n\times2n$, and they differ in which
matrix they take, in its sign, in the order of its rows and columns, and in the definition of the
Gaussian class. These notes make one choice throughout: the real antisymmetric Majorana covariance
$\Gamma$ in block order, with the sign for which the vacuum is $-\Omega_n$, and the Gaussian class
defined by Wick's rule. The paragraphs below translate each choice into the others and name the
statement where the translation is proved, or say that it is a computation on paper.}

\medskip
\noindent\revs{\textbf{The matrix.} The second moments of a density operator $\rho$ form the
two-point matrix $\G(\rho)_{ab}=\tr(\rho\,c_ac_b)=\delta_{ab}-\ii\,\Gamma(\rho)_{ab}$
(Section~\ref{sec:wick}), which part of the literature calls the correlation matrix\revs{, while
$\Gamma$ itself is often called the covariance matrix}. Its
symmetric part is fixed by the anticommutation relations, so $\Gamma$ carries the same information
in real form. The matrix of the annihilation operators $a_j=\tfrac12(c_{x_j}+\ii c_{y_j})$ follows
from $\G$ by expanding the products, a computation on paper:
\[
 \tr(\rho\,a_j^*a_k)=\tfrac12\delta_{jk}+\tfrac14\bigl(\Gamma_{x_jy_k}+\Gamma_{x_ky_j}\bigr)
 -\tfrac{\ii}4\bigl(\Gamma_{x_jx_k}+\Gamma_{y_jy_k}\bigr),
\]
whose diagonal $j=k$ is the occupation $\tfrac12(1+\Gamma_{x_jy_j})$ of mode $j$
(Section~\ref{sec:fock}).}

\medskip
\noindent\revs{\textbf{The sign.} The vacuum has $\Gamma=-\Omega_n$ (Corollary~\ref{cor:vacuum}),
and $\Gamma_{x_jy_j}$ is $-1$ for an empty and $+1$ for an occupied mode. The sign comes from the
factor $+\tfrac{\ii}{2}$ in Definition~\ref{def:density} together with the choice
$c_{y_j}=\ii Z_{<j}X_jZ_j$ of Definition~\ref{def:majorana}, under which $a_j$ annihilates the
vacuum. Replacing the factor by $-\tfrac{\ii}{2}$ replaces $\Gamma$ by $-\Gamma$, so that the
vacuum has $+\Omega_n$, the convention of part of the literature, and the ground covariance
$\ii\,\sgn(\ii h)$ of $H_h$ (after Theorem~\ref{thm:ground-state}) becomes $-\ii\,\sgn(\ii h)$.
Replacing $c_{y_j}$ by $-c_{y_j}$ conjugates $\Gamma$ by the reflection
$\operatorname{diag}(1_n,-1_n)$ of the $y$-labels; then $a_j$ annihilates the state with every mode
occupied, whose covariance is again $-\Omega_n$, and with $H_h$ written in the new operators the
ground covariance remains $\ii\,\sgn(\ii h)$, by a computation on paper. The
sign of the hopping in the Kitaev chain is a gauge (Remark~\ref{rem:kitaev-convention}).}

\medskip
\noindent\revs{\textbf{Order of the labels.} In the interleaved numbering of~\cite{bravyi}
and~\cite{ksvc}, $c_{2j-1}=x_j$ and $c_{2j}=y_j$. A matrix in block order becomes the matrix in
interleaved order by permuting rows and columns by $x_j\mapsto2j-1$, $y_j\mapsto2j$, which leaves
admissibility and purity unchanged (Lemma~\ref{lem:relabel}). In interleaved order $\Omega_n$
becomes $\bigoplus_j\begin{psmallmatrix}0&1\\-1&0\end{psmallmatrix}$ and the canonical form
$\D(\varepsilon)$ of Definition~\ref{def:block} becomes
$\bigoplus_j\varepsilon_j\begin{psmallmatrix}0&1\\-1&0\end{psmallmatrix}$.}

\medskip
\noindent\revs{\textbf{Gaussian unitaries.} A unitary $U$ implements a real matrix $O$ when
$Uc_aU^*=\sum_bO_{ab}c_b$; then $O$ is orthogonal and $\Gamma(U\rho U^*)=O^{\T}\Gamma(\rho)O$
(Theorem~\ref{thm:congruence}). A text that writes the Heisenberg form $U^*c_aU=\sum_bO'_{ab}c_b$
has $O'=O^{\T}$, and the same rule reads $\Gamma(U\rho U^*)=O'\,\Gamma(\rho)O'^{\T}$. The evolution
$e^{-\ii tH_h}$ implements $e^{-th}$ and moves covariances by $\Gamma\mapsto e^{th}\Gamma e^{-th}$
(Theorem~\ref{thm:evolution}).}

\medskip
\noindent\revs{\textbf{Contraction.} The contraction
$\Gamma\star\Gamma_p=A+B\,\Gamma_p(1-D\Gamma_p)^{-1}C$ of Definition~\ref{def:contract} is written
for any pair covariance $\Gamma_p$. On Fock space, for a Gaussian state $\rho$ with covariance
$\Gamma$, a Gaussian state $\sigma$ on the eliminated modes with covariance $\Gamma_p$, and
$1-D\Gamma_p$ invertible, it is the covariance of the normalized state $\rho$ traced against
$\sigma$ (Corollary~\ref{cor:contract-state} and Theorem~\ref{thm:normalization}). For antisymmetric $\Gamma$ and pure $\Gamma_p$, as for the bonds of a
chain, it equals $A+B(D+\Gamma_p)^{-1}B^{\T}$ (Theorem~\ref{thm:pure-bond}), the form
of~\cite[Sec.~II]{ksvc}. The sign of that formula depends on the orientation of the pair, which
Definition~\ref{def:bond} fixes and Theorem~\ref{thm:teleport} checks; reversing the orientation
replaces $\Gamma_p$ by $-\Gamma_p$. On Fock space the bond is the state
$(|00\rangle+\ii|11\rangle)/\sqrt2$ (Theorem~\ref{thm:bondstate}); it differs from
$|00\rangle+|11\rangle$ by a phase on one mode, and that phase makes the bond pair $x$ with $x$ and
$y$ with $y$.}

\medskip
\noindent\revs{\textbf{The Gaussian class.} A density operator is Gaussian here when its moments
are given by Wick's rule from its own two-point matrix (Definition~\ref{def:wick}); Wick's rule is
a condition beyond the second moments (Remark~\ref{rem:wick-condition}), \revf{and its recursion is the first-row expansion of
the Pfaffian form of~\cite{bravyi} (after Definition~\ref{def:wick})}. Gaussian states are often defined instead as normalized exponentials of
quadratic Hamiltonians together with their limits. The notes connect the two descriptions: the
Gibbs and thermal states are Gaussian (Corollary~\ref{cor:families}), every Gaussian state is a
rotated product mixture $U^*\rho_\lambda U$, a pure one is $U^*|s\rangle\langle s|U$, and one with
strictly admissible covariance is a Gibbs state (Corollary~\ref{cor:converse}). A Gaussian state is
determined by its covariance (Theorem~\ref{thm:unique}).}

\medskip
\noindent\revs{\textbf{Momentum space.} Section~\ref{sec:fourier} uses the characters
$\chi_k(i)=e^{-2\pi\ii\,ki/N}$ and expands the chain covariance in
$e^{+2\pi\ii\,k(i-j)/N}\,\Sigma(k)$; in the opposite convention the block at momentum $k$ is
$\Sigma(-k)=\overline{\Sigma(k)}$ (Theorem~\ref{thm:blocks}).}
